% =============================================================================
%  OP3: Operadic universality of interaction topology — skeleton paper
%  STATUS: UNPROVEN. Shares machinery with N1.
% =============================================================================
\documentclass[11pt]{article}

\usepackage[utf8]{inputenc}
\usepackage[T1]{fontenc}
\usepackage[expansion=false]{microtype}
\usepackage{amsmath, amssymb, amsthm, mathtools}
\usepackage[margin=1in]{geometry}
\usepackage{enumitem}
\usepackage[hidelinks]{hyperref}

\theoremstyle{plain}
\newtheorem{theorem}{Theorem}
\newtheorem*{namedopenproblem}{Open Problem}
\theoremstyle{definition}
\newtheorem{definition}[theorem]{Definition}

\title{Operadic Universality of Multiplex Interaction Topology \\
       \large (UNPROVEN --- skeleton paper for OP3)}
\author{Liam McGuire\thanks{Unaffiliated researcher, Seattle, USA. liammcg44@gmail.com}}
\date{\today}

\begin{document}
\maketitle

\begin{abstract}
This paper is a \textbf{skeleton} that records OP3 from the
scaffolding paper~\cite{emergent_systems_2026}. \textbf{No proof is
attempted.} OP3 shares machinery with N1 (both use the operad of
wiring diagrams) and should be developed jointly. The result is
UNPROVEN.
\end{abstract}

\section{Statement (UNPROVEN)}

\begin{quote}
\emph{``OP3 (interaction-topology expressiveness): is the
multiplex-with-$\diamond_\rho$ formalism universal in the operad of
wiring diagrams~\cite{vagner2015algebras}, in the sense that every
countably-many-layer multi-coupling-regime variation kernel factors
through a wiring-diagram presentation?''}
\end{quote}
\noindent (Scaffolding paper~\cite{emergent_systems_2026},
Conjectures section.)

\section{Plan of attack (no proof attempted)}

A two-stage Vagner--Spivak--Lerman / Yau-style proof, following the
sketch in \texttt{proof\_techniques.md} §OP3.

\begin{enumerate}[leftmargin=2em,itemsep=4pt]
  \item \textbf{Define $\mathcal{O}_W^{\mathrm{multi}}$ (UNPROVEN).}
    A coloured symmetric operad whose colours are pairs $(X,
    \mathcal{M}_+(X))$, operations are weighted hyperedge wiring
    diagrams with layer labels, composition is concatenation modulo
    associativity.
  \item \textbf{$W$-algebra on $V_T$ (UNPROVEN).} The candidate
    algebra sends each colour to the set of variation kernels and
    each operation to the $V_T$ construction implemented in the
    scaffolding paper §\texttt{sec:vt}.
  \item \textbf{Generator factorisation (UNPROVEN).} Every
    multiplex-with-$\diamond_\rho$ kernel factors through Yau's
    generators (identity, single hyperedge, layer-mixture,
    $\rho$-fusion). Verify by case analysis.
  \item \textbf{Scope conditions (UNPROVEN).} State conditions
    explicitly: countable layer index, measurable coupling rules,
    symmetric $\rho$, no higher-order coupling. Some of these are
    natural restrictions; others may need a non-symmetric operad
    upgrade.
\end{enumerate}

\section{Coordination with N1}

OP3 and N1 share the operadic machinery. The Decomp functor of N1
should be constructed on $\mathcal{O}_W^{\mathrm{multi}}$ rather than
on plain $\mathcal{O}_W$. Develop the two papers in lockstep.

\section{Status and follow-up}

\textbf{Current status: UNPROVEN.} Develop alongside
`n1_necessity/`.

\bibliographystyle{plain}
\bibliography{references}

\end{document}
